% Part: lambda-calculus
% Chapter: introduction
% Section: minimization

\documentclass[../../../include/open-logic-section]{subfiles}

\begin{document}

\olfileid{lam}{int}{min}
\olsection{\usetoken{S}{lambda definable} অপেক্ষকগুলি ন্যূনীকরণের অধীনে বদ্ধ}

\begin{lem}
ধরা যাক, $f(x,y)$ !!{lambda definable}। $g$-কে সংজ্ঞায়িত করা হলো
\[
g(x) \simeq \umin{y}{f(x,y)}.
\]
তাহলে $g$ !!{lambda definable}।
\end{lem}

\begin{proof}
মোটামুটি ধারণাটি হলো এই। $x$ দেওয়া থাকলে স্থির-বিন্দু ল্যাম্বডা পদ~$Y$
ব্যবহার করে এমন একটি অপেক্ষক $h_x(n)$ সংজ্ঞায়িত করব, যা~$n$ থেকে শুরু
করে একটি~$y$ খোঁজে; তখন $g(x)$ হলো ঠিক $h_x(0)$। অপেক্ষক~$h_x$-কে একটি
স্থির-বিন্দু সমীকরণের সমাধান হিসেবে প্রকাশ করা যায়:
\[
h_x(n) \simeq
\begin{cases}
n & \text{যদি $f(x,n) = 0$ হয়} \\
h_x(n+1) & \text{অন্যথায়।}
\end{cases}
\]

এবার বিস্তারিত দিই। যেহেতু $f$ ল্যাম্বডা-সংজ্ঞেয়, তাই কোনো পদ~$F$ তাকে
!!{lambda defined} করে।
% BN-SRC-274 TARGET-CORRECTION-BEGIN
\par\smallskip\noindent\textnormal{\small\textbf{উৎস-সংশোধন (BN-SRC-274):} লেমাটির অনুমান কেবল~$f$-এর ল্যাম্বডা-সংজ্ঞেয়তা, কিন্তু হিমায়িত প্রমাণটি অকারণে বলে যে~$f$ আদিম পুনরাবৃত্ত। বাংলা প্রমাণে ঘোষিত অনুমানটিই ব্যবহার করা হয়েছে: ল্যাম্বডা-সংজ্ঞেয়তা থেকে সংজ্ঞাকারী পদ~$F$-এর অস্তিত্ব সরাসরি মেলে।} % BN-SRC-274 TARGET-CORRECTION-END
মনে রেখো, আমাদের এমন একটি ল্যাম্বডা পদ~$D$-ও আছে, যাতে $D(M, N,
\bar{0}) \red M$ এবং $D(M, N, \bar{1}) \red N$। আপাতত~$x$ স্থির রেখে
$h_x$-কে !!{lambda define} করার জন্য এমন একটি পদ~$H$ খুঁজতে চাই, যা
$x$-এর উপর নির্ভরশীল এবং নিচের সমীকরণটি পূরণ করে:
\[
H(\num{n}) \equiv D(\num{n}, H(S(\num{n})), F(x, \num{n})).
\]
স্থির-বিন্দু পদ~$Y$ ব্যবহার করে এটি করা যায়। প্রথমে $U$-কে নিচের পদটি
ধরো:
\[
\lambd[h][\lambd[z][D(z,(h(Sz)),F(x,z))]],
\]
তারপর $H$-কে $YU$ পদটি ধরো। লক্ষ করো, $H$-এর একমাত্র মুক্ত চলরাশি হলো
$x$। এবার দেখাই, $H$ উপরের সমীকরণটি পূরণ করে।

$Y$-এর সংজ্ঞা থেকে পাই
\[
H = YU \equiv U(YU) = U(H).
\]
বিশেষত, প্রতিটি স্বাভাবিক সংখ্যা~$n$-এর জন্য
\begin{align*}
H(\num{n}) & \equiv U(H, \num{n}) \\
& \red D(\num{n}, H(S(\num{n})), F(x, \num{n})),
\end{align*}
যেমনটি দরকার ছিল। লক্ষ করো, শেষ পঙ্‌ক্তিতে~$x$-এর জায়গায় কোনো সংখ্যাপদ
$\num{m}$ বসালে, $F(\num{m}, \num{n})$ যদি $\num{0}$-তে হ্রাস পায় তবে
রাশিটি $\num{n}$-তে হ্রাস পায়; আর $F(\num{m}, \num{n})$ অন্য যেকোনো
সংখ্যাপদে হ্রাস পেলে রাশিটি $H(S(\num{n}))$-এ হ্রাস পায়।

প্রমাণ শেষ করতে $G$-কে $\lambd[x][H(\num 0)]$ ধরো। তখন $G$, $g$-কে
!!{lambda define}s; অর্থাৎ প্রতিটি~$m$-এর জন্য $g(m)$ সংজ্ঞায়িত হলে
$G(\num m)$, $\overline {g(m)}$-তে হ্রাস পায়, আর অন্যথায় তার কোনো
স্বাভাবিক রূপ নেই।
\end{proof}

\end{document}
